\section{A strict radial theorem for every positive Stieltjes transform}
\label{radpick:sec:main}
The preceding interlacing proof gives a nonnegative radial argument
derivative. A direct symmetrization strengthens it to strict positivity,
gives a quantitative bound, and identifies the sharp universal radius
range. This also removes the integer-leading restriction from angular
motion at depth two when combined with the earlier real-order factorization.

\begin{theorem}[Strict radial phase and its quantitative bound]
\label{radpick:thm:strict}
Let $\nu$ be a nonzero finite positive measure on $[0,1]$, and put
\[
 Q(z)=\int_{[0,1]}\frac{d\nu(u)}{1-uz},\qquad
 M_j=\int_{[0,1]}u^j\,d\nu(u).
\]
For $0<\theta<\pi$ and $0<\rho\le1$, the continuous phase satisfies
\begin{equation}\label{radpick:eq:quantitative}
 \partial_\rho\adoArg Q(\rho e^{\iu\theta})
 \ge \frac{4\rho^2\sin^3\theta(1-\cos\theta)^2}
              {(1+\rho)^8}\frac{M_1M_2}{M_0^2}.
\end{equation}
In particular, this derivative is strictly positive unless $\nu$ is
supported at zero, in which case $Q$ is constant. The range $\rho\le1$
is sharp as a universal assertion over positive measures on $[0,1]$:
at every prescribed radius greater than one a two-atom measure gives
a strictly negative radial phase derivative on the entire open upper arc.
\end{theorem}
\begin{proof}
Write $c=\cos\theta$, $h=\sin\theta$, $x=\rho u$, $y=\rho v$ and
$D_x=1-2cx+x^2$. These denominators are positive on the arc, and $Q$
is nonzero: its imaginary part is positive unless $\nu$ is supported at
zero, where it is a positive constant. Local domination justifies
differentiation and the double integral for
$\Imag(e^{\iu\theta}Q'(z)\overline{Q(z)})$.
Rationalizing and symmetrizing in $u,v$ gives
\begin{equation}\label{radpick:eq:pair}
 \partial_\rho\adoArg Q(\rho e^{\iu\theta})
 =\frac{h}{\rho|Q(\rho e^{\iu\theta})|^2}
   \int\!\!\int\frac{\mathcal P(x,y,c)}{D_x^2D_y^2}
                           \,d\nu(u)\,d\nu(v),
\end{equation}
where
\[
 \mathcal P=\frac12\{xD_y^2+yD_x^2
                    -(x-y)^2(xD_y+yD_x)\}.
\]
The exact polynomial identity that determines the sign is
\begin{align}\label{radpick:eq:positive-decomposition}
 \mathcal P={}&\frac{(1-x)(1-y)}2
   \{(x-y)^2+x(1-y)^2+y(1-x)^2\}\notag\\
 &+4xy(1-x)(1-y)(1-c)
       +2xy(x+y)(1-c)^2.
\end{align}
Every term is nonnegative for $0\le x,y\le1$ and $-1<c<1$.
The last term is strictly positive when $x,y>0$, including $x=y=1$.
Thus a positive mass on $(0,1]$ gives a strictly positive integral in
\eqref{radpick:eq:pair}; no approximation argument is needed for strictness.

For the explicit lower bound, retain only the last term and use
$D_x,D_y\le(1+\rho)^2$. Then
\[
 \partial_\rho\adoArg Q
 \ge\frac{4\rho^2h(1-c)^2M_1M_2}
              {(1+\rho)^8|Q|^2},
\]
since $\int\!\!\int uv(u+v)\,d\nu(u)\,d\nu(v)=2M_1M_2$.
Also $|1-te^{\iu\theta}|^2=(t-c)^2+h^2\ge h^2$, so
$|Q|\le M_0/h$. Substitution proves
\eqref{radpick:eq:quantitative}. All operations also apply at $\rho=1$,
because the chosen arc point is away from the cut.

For sharpness, let $R>1$, choose $R^{-2}<v<1$, and take
$\nu=v\delta_0+(1-v)\delta_1$. Its transform is
\[
 Q_v(z)=\frac{1-vz}{1-z}.
\]
Direct differentiation gives
\begin{equation}\label{radpick:eq:outside-counterexample}
 \partial_\rho\adoArg Q_v(\rho e^{\iu\theta})
 =\frac{(1-v)(1-\rho^2v)\sin\theta}
 {(1-2\rho\cos\theta+\rho^2)
  (1-2\rho v\cos\theta+\rho^2v^2)}.
\end{equation}
It is strictly negative at $\rho=R$ for every $0<\theta<\pi$.
This proves sharpness of the uniform disk range, without asserting
that every individual transform loses monotonicity just outside it.
\end{proof}

\begin{corollary}[All-real-order inward angular motion at depth two]
\label{radpick:cor:real-orders}
For every $a\ge0,b>0$, let $\theta_{a,b}(\rho)$ be the unique
upper-semicircle imaginary zero of the principal continuation of $F_{a,b}$.
Then
\[
 \theta_{a,b}'(\rho)<0,\qquad 0<\rho\le1.
\]
The derivative at one is the derivative of the local analytic branch,
equivalently its left derivative. For $0<\rho<\sigma\le1$ one also has
\[
 \frac{\sigma\cos\theta_{a,b}(\sigma)}
      {\rho\cos\theta_{a,b}(\rho)}>\frac\sigma\rho,
\]
so the Cartesian abscissa grows faster than the radius in this precise
ratio sense. This is distinct from the sign of
$\partial_\rho(\cos\theta/\rho)$.
\end{corollary}
\begin{proof}
Theorems~\ref{subunit:thm:zerofree} and \ref{subpick:thm:global} supply,
throughout the parameter domain, a factorization
\[
 F_{a,b}(z)=\frac{z^2Q_{a,b}(z)}{1-rz},\qquad 0<r\le1,
\]
with $Q_{a,b}$ a finite positive Stieltjes transform. In the
supercritical region with $a>0$, $r$ is the interior density crossing;
on the critical/subcritical region and the entire $a=0$ axis one may
take $r=1$. The positive measures have mass $2^{-a}$ and positive
density except possibly at the crossing. For an upper-arc point their
continuous unwrapped phase is
\[
 \Phi=2\theta-\adoArg(1-r\rho e^{\iu\theta})
                         +\adoArg Q_{a,b}(\rho e^{\iu\theta}).
\]
The resolvent sector bounds give $0<\Phi<2\pi$. Hence at the unique
imaginary zero, $\Phi=\pi$ and $F_{a,b}$ is strictly negative real.
The preceding real-order zero theorems prove
$\partial_\theta\Imag F_{a,b}<0$ there. Since
$\Imag F=|F|\sin\Phi$, this is equivalent to $\Phi_\theta>0$.
At fixed angle,
\[
 \Phi_\rho=\frac{r\sin\theta}
        {1-2r\rho\cos\theta+r^2\rho^2}
                 +\partial_\rho\adoArg Q_{a,b}>0.
\]
Implicit differentiation of $\Phi(\theta_{a,b}(\rho),\rho)=\pi$
therefore gives $\theta'=-\Phi_\rho/\Phi_\theta<0$.
The analytic continuation and angular simplicity justify this argument
at the endpoint and on parameter regions whose boundary series fail
the ordinary term test. Finally the unique angle lies in $(0,\pi/2)$,
so its strict decrease makes $\cos\theta$ strictly increase. The
displayed Cartesian ratio follows immediately.
\end{proof}

The direct pair identity and positive decomposition are independently
re-expanded symbolically and tested by exact rational arithmetic in
\src{verification/certify_stieltjes_radial.py}. Separate controls check
one-atom normalization and two-atom examples outside the disk, including
a prescribed-radius counterexample approaching radius one. These finite
checks catch algebraic mistakes; the measure theorem follows from the
sign decomposition and dominated integral argument, and the all-real-order
consequence uses the previously established positive factorization.
